\chapter{Cyclic $\Ainf$-Structures}
\label{ch:cyclic-ainf}

A Calabi--Yau category enters this volume through a single structural datum: a cyclic $\Ainf$-algebra of dimension $d$. The correspondence programme $\{\Phi_d\}$ to chiral algebras, the chiral modular characteristic $\kappa_{\mathrm{ch}}$, and the algebraisation-dependent kappas of the $\kappa_\bullet$-spectrum all read this input. This chapter fixes the definitions, records the examples at $d \in \{1, 2, 3\}$, and states the bridge to Part~\ref{part:bridge} precisely. The content is classical (Stasheff, Kontsevich, Keller, Costello); the Vol~III role is the identification of the cyclic dimension $d$ with the CY dimension appearing in the CY-A programme.

\begin{remark}[Cross-chapter dependencies]
\label{rem:cyclic-ainf-cross-deps}
The cyclic $\Ainf$-input fixed here is the data acted on by the correspondence programme $\{\Phi_d\}$ of Theorem~\textup{\ref{thm:phi-platonic}} (Chapter~\textup{\ref{ch:cy-to-chiral}}), characterised by the four universal properties \textup{(U1)}--\textup{(U4)}. The dimension $d$ of the cyclic structure is the integer entering the dimension stratification of $\kappa_{\mathrm{ch}}$ values in Theorem~\textup{\ref{thm:kappa-stratification-by-d}} (Chapter~\textup{\ref{ch:cy-d-kappa-stratification}}); for K3-fibered CY$_3$s, the Borcherds weight $\kappa_{\mathrm{BKM}} = c_N(0)/2$ of Proposition~\textup{\ref{prop:bkm-weight-universal}} supplies the second algebraisation invariant. On the Vol~I side, the cyclic bar complex of Section~\textup{\ref{sec:cyclic-structures}} is the chain-level engine of the universal Maurer--Cartan element of Vol~I Chapter~\textup{\ref{chap:climax-platonic}}, and the per-family Koszul conductor $K = -c_{\mathrm{ghost}}(\BRST)$ of Vol~I Chapter~\textup{\ref{chap:universal-conductor}} fixes the additive constants.
\end{remark}

\begin{remark}[Convention: $\kappa_\bullet$ subscripts]
\label{rem:cyclic-ainf-kappa-convention}
Every occurrence of $\kappa$ in Vol~III carries a subscript from the set $\{\mathrm{ch}, \mathrm{cat}, \mathrm{BKM}, \mathrm{fiber}\}$; when discussing the spectrum as a whole we write $\kappa_\bullet$. The four entries split into manifold invariants ($\kappa_{\mathrm{cat}}, \kappa_{\mathrm{fiber}}$), fixed by the geometry, and algebraisation invariants ($\kappa_{\mathrm{ch}}, \kappa_{\mathrm{BKM}}$), which depend on the construction applied to the manifold. Only the algebraisation pair is visible to the cyclic $\Ainf$-input through $\Phi$; the manifold pair is read off directly from the Hodge structure and the Mukai lattice.
\end{remark}

\section{$\Ainf$-algebras}
\label{sec:ainf-algebras}

The derived category of a CY manifold is not a strict associative algebra: the composition of $\Ext$-classes is associative only up to homotopy, and the homotopy itself is associative only up to a higher homotopy. Stasheff's $\Ainf$-formalism records the entire tower of homotopies as primary data $\mu_n$ rather than secondary corrections.

\begin{definition}[$\Ainf$-algebra]
\label{def:ainf-algebra}
An \emph{$\Ainf$-algebra} over a field $k$ is a $\Z$-graded $k$-vector space $A$ equipped with operations
\[
 \mu_n \colon A^{\otimes n} \longrightarrow A, \qquad \deg(\mu_n) = 2 - n, \quad n \geq 1,
\]
satisfying, for every $n \geq 1$, the $\Ainf$-relation
\[
 \sum_{\substack{p+q = n+1 \\ 1 \leq i \leq p}} (-1)^{\dagger(i,p,q)}\,
 \mu_p\bigl(a_1, \ldots, a_{i-1},\, \mu_q(a_i, \ldots, a_{i+q-1}),\, a_{i+q}, \ldots, a_n\bigr) = 0,
\]
where $\dagger(i,p,q) = (i-1)(q-1) + q \sum_{j=1}^{i-1}|a_j|$ follows the Koszul sign rule (see Appendix~\ref{app:conventions}).
\end{definition}

Specializing to small $n$ yields the familiar relations:
\begin{itemize}
 \item $n = 1$: $\mu_1^2 = 0$, so $(A, \mu_1)$ is a cochain complex.
 \item $n = 2$: $\mu_1(\mu_2(a,b)) = \mu_2(\mu_1 a, b) + (-1)^{|a|} \mu_2(a, \mu_1 b)$, so $\mu_1$ is a derivation of $\mu_2$.
 \item $n = 3$: $\mu_2$ is associative up to the explicit homotopy $\mu_3$, giving the Stasheff associahedron relation.
\end{itemize}
An $\Ainf$-algebra with $\mu_n = 0$ for $n \geq 3$ is a differential graded associative algebra. The $\mu_n$ for $n \geq 3$ record the higher homotopies that obstruct strict associativity after homological projection (Kadeishvili's theorem: any dg algebra transfers to a minimal $\Ainf$-algebra on its cohomology).

The $\Ainf$-formalism originates with Stasheff's associahedra~\cite{Stasheff1963}. Its modern use, as the natural habitat for derived categories, mirror symmetry, and deformation quantization, is due to Kontsevich and Soibelman~\cite{KontsevichSoibelman2009} and Keller~\cite{Keller2001Ainfty}. For a CY category $\cC$, the derived category $D^b(\cC)$ lifts canonically to an $\Ainf$-enhancement; this enhancement is unique up to quasi-isomorphism (Lunts--Orlov, Canonaco--Stellari). The cochain-level (non-minimal) model is the input used by the Vol~III functor: cyclicity is chain-level data and is not recoverable from the minimal model on cohomology alone.

\begin{remark}[$\Ainf$-category vs $\Ainf$-algebra]
\label{rem:ainf-cat-vs-algebra}
An $\Ainf$-category is the many-object generalization: objects $\mathrm{Ob}(\cC)$ and Hom-complexes $\Hom_\cC(X,Y)$ with composition maps $\mu_n \colon \Hom(X_0,X_1) \otimes \cdots \otimes \Hom(X_{n-1},X_n) \to \Hom(X_0,X_n)[2-n]$. An $\Ainf$-algebra is the one-object case. In this chapter we state definitions for algebras; the category version is verbatim by tagging arguments with source and target objects.
\end{remark}

\section{Cyclic structures and CY dimension}
\label{sec:cyclic-structures}

A strictly associative algebra equipped with an invariant non-degenerate pairing is a Frobenius algebra; the pairing satisfies $\langle ab, c \rangle = \langle a, bc \rangle$. Once the multiplication is replaced by a tower $\{\mu_n\}$, a single cyclic identity for $\mu_2$ is no longer sufficient: each higher product must individually respect the pairing, and the resulting tower of identities is what carries the $\bS^d$-framing on Hochschild homology.

\begin{definition}[Cyclic $\Ainf$-algebra]
\label{def:cyclic-ainf-d}
A \emph{cyclic $\Ainf$-algebra of dimension $d$} is a pair $\bigl((A, \{\mu_n\}_{n \geq 1}),\, \langle -, - \rangle\bigr)$, where $(A, \{\mu_n\})$ is an $\Ainf$-algebra and
\[
 \langle -, - \rangle \colon A \otimes A \longrightarrow k[-d]
\]
is a non-degenerate graded symmetric pairing of degree $-d$ (equivalently, $\langle a, b \rangle$ is nonzero only when $|a| + |b| = d$), satisfying the \emph{cyclic invariance} identity
\[
 \langle \mu_n(a_1, a_2, \ldots, a_n), a_{n+1} \rangle
 \;=\; (-1)^{\epsilon_n}\,
 \langle a_1, \mu_n(a_2, \ldots, a_{n+1}) \rangle
 \qquad \text{for all } n \geq 1,
\]
with sign $\epsilon_n = n + |a_1|(|a_2| + \cdots + |a_{n+1}|)$ determined by the Koszul rule.
\end{definition}

The integer $d$ is the \emph{CY dimension} of the cyclic $\Ainf$-algebra. In the language of Costello~\cite{Costello2005TCFT,Costello2007Ainfty}, a cyclic $\Ainf$-algebra of dimension $d$ is precisely an algebra over the operad of framed chains on the moduli of disks with marked points, with a $d$-shifted symplectic structure. This operadic viewpoint is what produces, after homotopy transfer, the $\bS^d$-framing of Hochschild homology used in Chapter~\ref{ch:cy-to-chiral}.

\begin{proposition}[Cyclic invariance on the Hochschild complex]
\label{prop:cyclic-invariance-hochschild}
\ClaimStatusProvedElsewhere{}
Let $(A, \langle -, - \rangle)$ be a cyclic $\Ainf$-algebra of dimension $d$. The pairing induces a non-degenerate pairing on the Hochschild complex
\[
 \langle -, - \rangle_{\mathrm{HH}} \colon \HH_\bullet(A) \otimes \HH_{d-\bullet}(A) \longrightarrow k
\]
implementing Poincar\'e duality of dimension $d$. Equivalently, the cyclic structure lifts Hochschild homology to negative cyclic homology with a class in $\HC^-_d(A)$ that is closed under the Connes differential.
\end{proposition}

\begin{proof}[Attribution]
Costello~\cite{Costello2007Ainfty} Thm.~A; Kontsevich--Soibelman~\cite{KontsevichSoibelman2009} Sec.~10. The distinction is essential: the CY trace lives in $\HC^-_d$, not in $\HH_d \to k$, because the $S^1$-invariance of the cyclic pairing is needed to produce the $\bS^d$-framing.
\end{proof}

\begin{remark}[Geometric origin of $d$]
\label{rem:d-geometric}
For $X$ a smooth projective CY $n$-fold, the derived category $D^b(\mathrm{Coh}(X))$ carries a cyclic $\Ainf$-enhancement of dimension $d = n$. The pairing is the Serre-duality pairing
\[
 \langle -, - \rangle_{\mathrm{Serre}} \colon
 \Ext^i(\cF, \cG) \otimes \Ext^{n-i}(\cG, \cF) \longrightarrow \Ext^n(\cF, \cF) \xrightarrow{\mathrm{tr}} H^n(X, \omega_X) \cong k,
\]
using the trivialization $\omega_X \cong \cO_X$ that defines the CY structure. $d$ is the complex dimension, not the real dimension $2n$. The Serre functor $\mathbb{S} = (-) \otimes \omega_X[n]$ becomes $[n]$-shift under the CY trivialization, which is what produces the $d$-shifted pairing.
\end{remark}

\begin{remark}[Cyclic bar complex]
\label{rem:cyclic-bar-complex}
The cyclic pairing enters the bar complex $B(A) = T^c(s^{-1}\bar A)$ through the cyclic quotient $\mathrm{CC}_\bullet(A) = B(A)/(1 - t)$ where $t$ is the signed cyclic rotation. The factor $s^{-1}$ desuspends: $|s^{-1}v| = |v| - 1$. The augmentation ideal $\bar A = \ker(\varepsilon)$ is used rather than $A$ itself. The cyclic bar complex is the primary invariant of $(A, \mu_n, \langle-,-\rangle)$ and is what Part~\ref{part:bridge} promotes to a factorization coalgebra on curves.
\end{remark}

\begin{remark}[Cyclic symmetry at $n = 2$]
\label{rem:cyclic-n2}
Specializing the cyclic invariance identity to $n = 2$ gives
\[
 \langle \mu_2(a_1, a_2), a_3 \rangle \;=\; (-1)^{\epsilon_2}\, \langle a_1, \mu_2(a_2, a_3) \rangle,
\]
which for a strictly associative cyclic algebra ($\mu_n = 0$ for $n \geq 3$) reduces to the standard invariance condition for a Frobenius algebra of degree $d$. The higher-$n$ cases are the $\Ainf$-correction: $\mu_3, \mu_4, \ldots$ must individually respect the pairing. In the $n = 3$ case, the identity
\[
 \langle \mu_3(a_1, a_2, a_3), a_4 \rangle \;=\; (-1)^{\epsilon_3}\, \langle a_1, \mu_3(a_2, a_3, a_4) \rangle
\]
is the first genuinely homotopical constraint and is what fails in a naive restriction of the Fukaya category to a sub-Lagrangian: cyclic invariance distinguishes the true CY input from merely $\Ainf$-enhanceable ones.
\end{remark}

\section{Examples}
\label{sec:cyclic-ainf-examples}

The examples below realise the three CY dimensions $d \in \{1, 2, 3\}$ relevant to Vol~III. Each fixes one numerical invariant visible to the correspondence $\Phi_d$.

\begin{example}[$d = 1$: $\mathfrak{sl}_2$ as cyclic $\Ainf$-algebra]
\label{ex:cy1-sl2}
Let $\mathfrak{g}$ be a finite-dimensional simple Lie algebra with non-degenerate invariant form $\langle -, - \rangle_{\mathfrak{g}}$. The Chevalley--Eilenberg cochain complex $C^\bullet(\mathfrak{g})$ is a dg commutative algebra of dimension $\dim \mathfrak{g}$; applying Koszul duality twice produces a cyclic $\Ainf$-algebra of CY dimension $1$ whose underlying vector space is $\mathfrak{g}[1]$ with pairing $\langle -, - \rangle_{\mathfrak{g}}$ (shifted to degree $-1$). For $\mathfrak{g} = \mathfrak{sl}_2$, fix the basis $\{e[1], f[1], h[1]\}$ of $\mathfrak{sl}_2[1]$ in degree $-1$. The binary operation $\mu_2$ is the desuspended Lie bracket:
\[
 \mu_2(e[1], f[1]) = h[1], \qquad \mu_2(h[1], e[1]) = 2\, e[1], \qquad \mu_2(h[1], f[1]) = -2\, f[1],
\]
with the other values determined by graded antisymmetry. The cyclic pairing has nonzero values
\[
 \langle e[1], f[1] \rangle = 1, \qquad \langle h[1], h[1] \rangle = 2,
\]
which is the trace form $\mathrm{tr}_{\mathrm{def}}(XY)$ on the defining representation, shifted to degree $-1$, producing a CY structure of dimension $d = 1$. (The Killing form $B(h,h) = 8$ differs by a factor of $4 = 2 h^\vee$.) The $\Ainf$-relations reduce to the Jacobi identity, and all $\mu_n$ for $n \geq 3$ vanish. Under $\Phi$ this produces the affine Kac--Moody algebra $\widehat{\mathfrak{sl}}_2$ at the critical level: the simplest nontrivial entry of the CY landscape.
\end{example}

\begin{example}[$d = 1$ from an elliptic curve]
\label{ex:cy1-elliptic}
Let $E$ be an elliptic curve over $\C$. The derived category $D^b(\mathrm{Coh}(E))$ has a minimal cyclic $\Ainf$-enhancement of dimension $d = 1$, computed explicitly by Polishchuk~\cite{Polishchuk2011}. The generators are $\cO_E$ and a single degree-$1$ class in $\mathrm{Ext}^1(\cO_E, \cO_E) = H^1(E, \cO_E) \cong \C$. The Serre pairing $\mathrm{Ext}^1(\cO_E, \cO_E) \otimes \Hom(\cO_E, \cO_E) \to \C$ realizes the $d = 1$ cyclic structure. For higher-degree line bundles, Polishchuk's theta-function computation gives explicit formulas for $\mu_n$, $n \geq 3$, in terms of modular forms on $\Gamma_1(N)$. This example sits at the boundary between $d = 1$ (where the $\bS^1$-framing is classical) and the modular-form territory that Vol~I controls.
\end{example}

\begin{example}[$d = 2$ from K3]
\label{ex:cy2-k3}
Let $X$ be a smooth projective K3 surface. The derived category $D^b(\mathrm{Coh}(X))$ is a cyclic $\Ainf$-category of dimension $d = 2$ (Caldararu~\cite{Caldararu2005}). The Hochschild cohomology is
\[
 \HH^\bullet(D^b(\mathrm{Coh}(X))) \;\cong\; H^\bullet(X, \wedge^\bullet T_X) \;\cong\; H^\bullet(X, \C),
\]
using the CY trivialization $\wedge^2 T_X \cong \cO_X$. The Kontsevich HKR decomposition gives $\HH^0 \simeq k$, $\HH^1 \simeq 0$, $\HH^2 \simeq k^{22}$, $\HH^3 \simeq 0$, $\HH^4 \simeq k$ (Example~\ref{ex:hh-k3} of Chapter~\ref{ch:cy-categories}), so
\[
 \dim \HH^\bullet(K3) = 1 + 0 + 22 + 0 + 1 = 24,
\]
recovering the Euler characteristic $\chi(X) = 24$ familiar from lattice K3 geometry. The Mukai row-grading $\HH_{p,q}$ with $p = q$ collapses the same total to $(1+1) + 20 + (1+1) = 2 + 20 + 2 = 24$. The holomorphic Euler characteristic, which is the relevant invariant for the CY kappa-spectrum, is
\[
 \chi(\cO_X) \;=\; 2.
\]
This is the value $\kappa_{\mathrm{cat}}(\mathrm{K3}) = 2$ entering the Vol~III kappa-table.
\end{example}

\begin{example}[$d = 3$ from the quintic]
\label{ex:cy3-quintic}
Let $Q \subset \P^4$ be a smooth quintic threefold. $D^b(\mathrm{Coh}(Q))$ is a cyclic $\Ainf$-category of dimension $d = 3$. By homological mirror symmetry (Kontsevich's conjecture, established for the quintic at genus $0$ by Sheridan~\cite{Sheridan2015}), there is an $\Ainf$-equivalence
\[
 D^b(\mathrm{Coh}(Q)) \;\simeq\; D^\pi \mathrm{Fuk}(Q^\vee),
\]
where $Q^\vee$ is the mirror quintic and $\mathrm{Fuk}$ is the Fukaya category (wrapped, compact, depending on geometric input). Both sides are cyclic $\Ainf$ of dimension $d = 3$, and HMS is an equivalence of cyclic $\Ainf$-categories: the Serre pairing on the B-side matches the Poincar\'e pairing on Lagrangian intersections on the A-side. The Hochschild cohomology
\[
 \HH^\bullet(Q) \cong H^\bullet(Q, \wedge^\bullet T_Q)
\]
gives, using the CY$_3$ isomorphism $\wedge^r T_Q \simeq \Omega^{3-r}_Q$,
\[
 \HH^2(Q) \simeq H^1(Q, T_Q) \cong k^{101}
 \quad \text{(Kodaira--Spencer, $q+r = 2$ with $q=1,r=1$)},
\]
\[
 \HH^3(Q) \simeq H^0(\wedge^3 T_Q) \oplus H^1(\wedge^2 T_Q) \oplus H^2(T_Q) \oplus H^3(\cO_Q)
 \cong k^4
\]
(the summands identify with $H^0(\cO_Q), H^1(\Omega^1_Q), H^2(\Omega^2_Q), H^3(\cO_Q)$, of dimensions $1,1,1,1$, including the Yukawa-coupling generator $H^0(\wedge^3 T_Q) \cong H^0(\cO_Q)$). The holomorphic Euler characteristic is $\chi(\cO_Q) = 0$ by Serre duality on a CY$_3$, so the naive $\kappa_{\mathrm{cat}}$ vanishes; the nontrivial chiral data enters through the Kodaira--Spencer summand of $\HH^2$, not through the top pairing. The chiral algebra $A_Q = \Phi_3(D^b(\Coh(Q)))$ is constructed by Theorem~CY-A$_3$ (Theorem~\ref{thm:cy-to-chiral-d3}).
\end{example}

\section{The bridge to the Vol~III functor}
\label{sec:bridge-to-functor}

The cyclic $\Ainf$-structure is the complete input to the CY-to-chiral correspondence programme $\{\Phi_d\}_{d \geq 1}$. The $d$-th signature is
\[
 \Phi_d \colon \mathrm{CY}_d\text{-}\Cat \longrightarrow \mathrm{E}_{n(d)}\text{-}\mathrm{ChirAlg},
 \qquad (\cC, \mu_n, \langle-,-\rangle) \longmapsto A_\cC,
\]
where the operadic level $n(d) = n_{\mathrm{native}}(d) \in \{\infty, 2, 1, 1\}$ for $d \in \{1, 2, 3, \geq 4\}$ is dictated by property~\textup{(U3)} of Theorem~\textup{\ref{thm:phi-platonic}} (Gerstenhaber-bracket degree $1 - d$). At $d \geq 3$ the native output is $\Eone$; the braided $\Etwo$-structure is recovered on the Drinfeld centre $\cZ(\Rep^{\Eone}(A_\cC))$ via half-braidings. The centre is the right adjoint to the forgetful functor $\mathrm{BrMon} \to \mathrm{Mon}$, a categorified commutant construction, not a categorified averaging map (averaging projects $\Eone$ to $E_\infty$ and kills the Hopf structure). The data of $(\cC, \mu_n, \langle-,-\rangle)$ factors through $\Phi$ by the four-step passage in Chapter~\ref{ch:cy-to-chiral}: \emph{categorical} Hochschild cohomology with Gerstenhaber bracket (distinct from the chiral and topological Hochschild theories of the output and of Vol~II respectively, see Remark~\ref{rem:cyclic-ainf-three-hochschild}), Lie conformal algebra via the cyclic pairing, factorisation envelope, and $\Etwo$-enhancement via the $\bS^d$-framing. The cyclic pairing $\langle-,-\rangle$ is what becomes the invariant form on $A_\cC$; without it, the factorisation envelope is not a chiral algebra with a trace.

\begin{remark}[Three Hochschild theories enter here separately]
\label{rem:cyclic-ainf-three-hochschild}
The first step of $\Phi$ extracts the \emph{categorical} Hochschild cohomology $\HH^\bullet_{\mathrm{cat}}(\cC)$ of the cyclic $\Ainf$-input (Chapter~\textup{\ref{ch:hochschild-calculus}}). The output of $\Phi$ has its own \emph{chiral} Hochschild cohomology $\ChirHoch^\bullet(A_\cC)$, concentrated in cohomological degrees $\{0,1,2\}$ by Vol~I Theorem~H. These are NOT the same object: the categorical input is an $E_2$-algebra (Deligne; Tamarkin); the chiral output has the much sharper $\{0,1,2\}$ amplitude. The third Hochschild theory, \emph{topological} Hochschild $\HH^\bullet_{\mathrm{top}}$ for $E_1$-algebras (Vol~II's slab algebra of 3d HT theories), is a third functor on a third source. The three theories differ in their ambient category ($\Ainf$-categories, chiral algebras on a curve, $E_1$-algebras of spectra), in their derivation ($\RHom$ of bimodules, chiral endomorphisms on $\mathrm{Ran}$-space, $\THH$ of an $E_1$-ring), and in their cohomological concentration; Vol~III statements that attribute Hochschild data to $\cC$ always mean $\HH^\bullet_{\mathrm{cat}}$ unless another subscript is present.
\end{remark}

\begin{theorem}[Cyclic $\Ainf$ input determines $\kappa_{\mathrm{cat}}$ at $d = 2$, $h^{1,0} = 0$]
\label{thm:cyclic-ainf-kappa-cat-d2}
\ClaimStatusProvedHere{Restricted to $h^{1,0}(X) = 0$ (equivalently $\HH_{-1}(\cC) = 0$). The identification $\kappa_{\mathrm{cat}} = \kappa_{\mathrm{ch}}$ via the Hodge supertrace does not extend to odd $d$ (where $\chi(\cO_X) = 0$ by Serre duality but $\kappa_{\mathrm{ch}}$ is generically nonzero) or to $d = 2$ with $h^{1,0}(X) \neq 0$ (where the Beauville correction $h^{1,0}$ contributes). Counter-example in Remark~\textup{\ref{rem:cyclic-ainf-kappa-d2-hypothesis}}.}
Let $\cC$ be a smooth proper cyclic $\Ainf$-category of dimension $d = 2$ such that $\HH_{-1}(\cC) = 0$ (equivalently, in the geometric case $\cC = D^b(\Coh(X))$, the hypothesis $h^{1,0}(X) = 0$), and let $A_\cC = \Phi_2(\cC)$ be its image under the $d = 2$ CY-to-chiral correspondence (Theorem~\ref{thm:cy-to-chiral}). Then the categorical kappa is the Hodge supertrace
\[
 \kappa_{\mathrm{cat}}(\cC) \;=\; \chi^{\mathrm{CY}}(\cC) \;=\; \chi(\cO_X) \;=\; \sum_{q=0}^{d} (-1)^q\, h^{0,q}(X),
\]
and this equals the chiral modular characteristic $\kappa_{\mathrm{ch}}(A_\cC)$ of the image chiral algebra (Proposition~\ref{prop:cy-kappa-d2}; the Serre duality $\mathbb{S}_\cC \simeq [2]$ kills the one-loop correction from quantisation Step~4). For $\cC = D^b(\mathrm{Coh}(\mathrm{K3}))$, both invariants evaluate to $\kappa_{\mathrm{cat}}(\mathrm{K3}) = \kappa_{\mathrm{ch}}(\Phi(D^b(\Coh(K3)))) = 2$. At $d \geq 3$ the identification $\kappa_{\mathrm{ch}} = \kappa_{\mathrm{cat}}$ FAILS (Conjecture~\ref{conj:cy-kappa-identification}).
\end{theorem}

\begin{proof}[Proof sketch]
The cyclic pairing of dimension $d = 2$ produces a degree-$(-2)$ symmetric form on the Hochschild complex; under HKR (Theorem~\textup{\ref{thm:hkr-smooth}}) its top-degree component is a linear map $\HH^2(\cC) \to k$ whose dimension (as a quotient of $\HH^\bullet$) is $\chi(\cO_X)$ in the geometric case. Property~\textup{(U1)} of Theorem~\textup{\ref{thm:phi-platonic}} sends this trace to the invariant form on $A_\cC$, and the hypothesis $\HH_{-1}(\cC) = 0$ kills the would-be one-loop correction from quantisation (Hodge-supertrace mechanism, Remark~\textup{\ref{rem:kappa-cat-from-chi}}). The resulting $\kappa_{\mathrm{cat}}$ coincides (up to $\bS^2$-framing) with the coefficient of the $E_2$-algebra central extension. For K3 the computation reduces to $\chi(\cO_{\mathrm{K3}}) = 2$; see Chapter~\ref{ch:cy-to-chiral} Section~\ref{sec:cy-chiral-functor} and the compute module \texttt{cy\_to\_chiral\_functor.py} for the full verification.
\end{proof}

\begin{remark}[Hypothesis $h^{1,0} = 0$ is essential]
\label{rem:cyclic-ainf-kappa-d2-hypothesis}
The simply-connected hypothesis is not cosmetic. The abelian surface $T^4$ has $h^{1,0}(T^4) = 2$, so $\chi(\cO_{T^4}) = 1 - 2 + 1 = 0$, but additivity of $\kappa_{\mathrm{ch}}$ under tensor products gives $\kappa_{\mathrm{ch}}(\Phi(T^4)) = \kappa_{\mathrm{ch}}(\Phi(E)) + \kappa_{\mathrm{ch}}(\Phi(E)) = 1 + 1 = 2$, contradicting $\kappa_{\mathrm{cat}} = 0$. The Beauville extension of Proposition~\textup{\ref{prop:beauville-kappa-formula}} restores $\kappa_{\mathrm{ch}}(\Phi(D^b(\Coh(X)))) = \chi(\cO_X) + h^{1,0}(X)$ for arbitrary smooth proper $\CY_2$, recovering Theorem~\textup{\ref{thm:cyclic-ainf-kappa-cat-d2}} in the $h^{1,0} = 0$ case.
\end{remark}

\begin{remark}[Independent verification path]
\label{rem:cyclic-ainf-kappa-d2-verification}
The proof of Theorem~\textup{\ref{thm:cyclic-ainf-kappa-cat-d2}} uses (i) the cyclic $\Ainf$ pairing on $\cC$ and (ii) the chiral envelope step of $\Phi$. Two disjoint verification paths converge on the same numerical output. The lattice path (Mukai 1984): the Mukai pairing on $H^\bullet(K3, \Z)$ has rank $24$, recovering the topological Euler characteristic $\chi(K3) = 24$, independently of any $\Ainf$-enhancement. The Hodge-supertrace path: $\chi(\cO_{K3}) = h^{0,0} - h^{0,1} + h^{0,2} = 1 - 0 + 1 = 2$, computed from the K3 Hodge diamond independently of HKR. The lattice rank $24$ and the holomorphic Euler characteristic $2$ are distinct invariants with disjoint derivations; conflating them (e.g. claiming $\kappa_{\mathrm{cat}}(K3) = 24$ because $\dim \HH^\bullet = 24$) is the failure mode this independent verification blocks.
\end{remark}

\begin{conjecture}[Cyclic $\Ainf$ input determines $\kappa_{\mathrm{cat}}$ at $d = 3$]
\label{conj:cyclic-ainf-kappa-cat-d3}
\ClaimStatusConjectured{}
Let $\cC$ be a smooth proper cyclic $\Ainf$-category of dimension $d = 3$. By Theorem~CY-A$_3$ (Theorem~\ref{thm:cy-to-chiral-d3}), the $d = 3$ CY-to-chiral correspondence $\Phi_3$ exists with chain-level $\bS^3$-framing. Then
\[
 \kappa_{\mathrm{cat}}(\cC) \;=\; \chi^{\mathrm{CY}}(\cC),
\]
and for $\cC = D^b(\mathrm{Coh}(Q))$ with $Q$ a smooth quintic threefold, this evaluates to a specific integer determined by the $\bS^3$-framing data.
\end{conjecture}

\begin{remark}[Scope: what cyclic $\Ainf$ does not determine]
\label{rem:cyclic-ainf-scope}
The cyclic $\Ainf$-structure determines $\kappa_{\mathrm{cat}}$ and the chiral kappa $\kappa_{\mathrm{ch}}$ via the correspondence $\Phi_d$, but it does \emph{not} determine the BKM kappa $\kappa_{\mathrm{BKM}}$ or the fiber kappa $\kappa_{\mathrm{fiber}}$ of the kappa-spectrum. Those kappas arise from additional data (lattice embedding for $\kappa_{\mathrm{fiber}}$, Borcherds product structure for $\kappa_{\mathrm{BKM}}$) that is not visible to the cyclic $\Ainf$-category alone. The four-way kappa-spectrum for K3 $\times E$ is $\{\kappa_{\mathrm{cat}}, \kappa_{\mathrm{ch}}, \kappa_{\mathrm{BKM}}, \kappa_{\mathrm{fiber}}\} = \{0, 3, 5, 24\}$ (where $\kappa_{\mathrm{cat}}(K3 \times E) = \chi(\cO_{K3}) \cdot \chi(\cO_E) = 2 \cdot 0 = 0$ by K\"unneth; the value $2 = \chi(\cO_{K3})$ is the fiber kappa $\kappa_{\mathrm{cat}}(K3)$, not the total-space invariant). Only $\kappa_{\mathrm{ch}}$ is extracted from the cyclic $\Ainf$-input via the correspondence $\Phi_d$; $\kappa_{\mathrm{cat}}$ is a manifold invariant independent of the algebraization.
\end{remark}

\begin{remark}[Comparison with prior work]
\label{rem:cyclic-ainf-prior-work}
Four source threads feed the construction used here. Stasheff~\cite{Stasheff1963} introduced the associahedra and the higher homotopies $\mu_n$. Kontsevich~\cite{Kontsevich1995} identified cyclic $\Ainf$-algebras with algebras over the operad of ribbon graphs, providing the link to moduli of curves with boundary. Costello~\cite{Costello2005TCFT,Costello2007Ainfty} proved that cyclic $\Ainf$-categories are equivalent to open topological conformal field theories and supplied the first rigorous construction of the associated chain-level trace. Kontsevich--Soibelman~\cite{KontsevichSoibelman2009} axiomatized the CY structure in terms of the negative cyclic class and gave the formalism used in Part~\ref{part:cy-categories}. Keller~\cite{Keller2001Ainfty} surveys the homological-algebra side. For explicit computations on projective varieties, Polishchuk~\cite{Polishchuk2011} computed the cyclic $\Ainf$-structure on elliptic curves and on their products, and Caldararu~\cite{Caldararu2005} set up the Hochschild calculus for smooth proper CY categories. The Vol~III role is the specific mapping of this input through the correspondence programme $\{\Phi_d\}$, producing chiral algebras whose modular characteristic can be computed and compared across the four kappas of the spectrum.
\end{remark}

\begin{remark}[Costello's operadic description]
\label{rem:costello-operadic}
Costello~\cite{Costello2005TCFT} proves that cyclic $\Ainf$-algebras of dimension $d$ are equivalent to algebras over the open topological conformal field theory operad with $d$-shifted framing. In this formulation, the $\bS^d$-framing used in Chapter~\ref{ch:cy-to-chiral} is visible as the action of the framing circle on the moduli of disks with boundary marked points. For $d = 2$ the framing is abelian ($\pi_1(SO(2)) = \Z$) and the envelope step produces an honest $\Etwo$-algebra. For $d = 3$ the framing is nonabelian ($\pi_1(SO(3)) = \Z/2$) and the envelope step is the source of the $\bS^3$-framing chain-level obstruction analysed in Chapter~\ref{ch:m3-b2-saga} and resolved in the $\infty$-categorical CY-A$_3$ proof of Theorem~\ref{thm:derived-framing-obstruction}.
\end{remark}
